\chapter*{Kurzreferat}
Diese Arbeit behandelt die Konzeptionierung, Entwicklung und den Aufbau einer Wortuhr mit erweiterten Informations- und Zusatzfunktionen. 
Ziel war die Entwicklung eines funktionsfähigen Prototyps, der die klassische wortbasierte Zeitdarstellung mit zusätzlichen informationsorientierten Funktionen kombiniert. 
Hierzu wurden elektronische, softwaretechnische und mechanische Komponenten zu einem Gesamtsystem zusammengeführt.
Im Rahmen der Arbeit erfolgte zunächst die Analyse bestehender Wortuhrkonzepte sowie die Definition funktionaler und technischer Anforderungen. 
Darauf aufbauend wurden die Hardwarearchitektur, die Spannungsversorgung, die LED-Ansteuerung sowie die erforderlichen Schnittstellen ausgelegt. 
Zusätzlich wurde eine modulare Softwarestruktur entwickelt. 
Diese ermöglicht die Verarbeitung externer Datenquellen, die Steuerung verschiedener Anzeigeelemente sowie die Bereitstellung netzwerkbasierter Funktionen.
Der entwickelte Prototyp stellt die Uhrzeit in sprachlicher Form dar und erweitert die Anzeige um zusätzliche Informationen. 
Hierzu zählen unter anderem Wetterdaten, Temperaturinformationen, Luftfeuchtigkeit, Mondphasen sowie sekundengenaue Anzeigeelemente. 
Darüber hinaus wurde eine Steuerung über mobile Endgeräte mittels Webserver implementiert.
Die Ergebnisse zeigen, dass die Integration zusätzlicher Funktionen in eine Wortuhr technisch umsetzbar ist, ohne die charakteristische Darstellungsform wesentlich zu beeinträchtigen. 
Die Arbeit bildet damit eine Grundlage für zukünftige Erweiterungen im Bereich vernetzter und interaktiver Anzeige- und Informationssysteme.

\chapter*{Abstract}
% <text hier>